Seja \(p\) um número primo. Uma pulga está na posição \(0\) da reta real. A cada minuto, a pulga tem três possibilidades: ficar em sua posição, ou se mover \(1\) para a esquerda ou para a direita. Após \(p-1\) minutos, ela deseja estar novamente em \(0\). Denote por \(f(p)\) o número de suas estratégias para fazer isso (por exemplo, \(f(3)=3\): ela pode ficar em \(0\) durante todo o tempo, ou ir para a esquerda e depois para a direita, ou ir para a direita e depois para a esquerda). Determine \(f(p) \pmod p\).
